\documentclass[12pt,a4paper]{article}

\usepackage[utf8]{inputenc}
\usepackage[T1]{fontenc}
\usepackage[brazilian]{babel}
\usepackage{lmodern}
\usepackage[a4paper,margin=2.3cm]{geometry}
\usepackage{microtype}
\usepackage{parskip}
\usepackage{enumitem}
\usepackage{booktabs}
\usepackage{tabularx}
\usepackage{array}
\usepackage{xcolor}
\usepackage{listings}
\usepackage{hyperref}
\usepackage{fancyhdr}
\usepackage{lastpage}
\usepackage[most]{tcolorbox}

\definecolor{NSSBlue}{HTML}{174A7E}
\definecolor{NSSLight}{HTML}{EEF5FB}
\definecolor{NSSGreen}{HTML}{2E7D5A}
\definecolor{NSSGray}{HTML}{5F6B76}
\definecolor{NSSRed}{HTML}{9A3412}

\hypersetup{
  colorlinks=true,
  linkcolor=NSSBlue,
  urlcolor=NSSBlue,
  citecolor=NSSBlue,
  pdftitle={NSS Front-end - Arquitetura V1 e Navegacao Geografica},
  pdfauthor={Nucleo de Situacao de Saude - UENF}
}

\pagestyle{fancy}
\fancyhf{}
\lhead{NSS Front-end}
\rhead{Arquitetura V1}
\cfoot{\thepage\ de \pageref{LastPage}}
\setlength{\headheight}{15pt}

\newcolumntype{Y}{>{\raggedright\arraybackslash}X}
\newcolumntype{L}[1]{>{\raggedright\arraybackslash}p{#1}}

\newtcolorbox{decision}[1][]{
  colback=NSSLight,
  colframe=NSSBlue,
  title=Decisao arquitetural,
  fonttitle=\bfseries,
  #1
}

\newtcolorbox{warningbox}[1][]{
  colback=orange!7,
  colframe=NSSRed,
  title=Atencao,
  fonttitle=\bfseries,
  #1
}

\lstdefinestyle{code}{
  basicstyle=\ttfamily\small,
  breaklines=true,
  columns=fullflexible,
  frame=single,
  rulecolor=\color{black!20},
  backgroundcolor=\color{black!2},
  showstringspaces=false,
  tabsize=2
}
\lstset{style=code}

\title{\textbf{Nucleo de Situacao de Saude -- UENF}\\[0.4cm]
\Large Front-end: Arquitetura V1, Navegacao Geografica e Escopo de Demonstracao}
\author{Documentacao para implementacao incremental}
\date{15 de setembro de 2026}

\begin{document}
\maketitle
\thispagestyle{empty}
\vfill

\begin{decision}
Esta V1 prioriza uma demonstracao visual forte e funcional, mantendo a arquitetura definitiva do NSS. O mapa deixa de ser um elemento opcional e passa a ser a principal forma de navegacao da interface. O front-end conversa somente com a API Java quando estiver em modo real; enquanto os endpoints epidemiologicos ainda nao estiverem disponiveis, a mesma interface opera com dados mockados pelo mesmo contrato tipado.
\end{decision}

\vfill
\newpage
\tableofcontents
\newpage

\section{Objetivo desta V1}

Este documento define a primeira versao arquitetural do repositorio \texttt{nss-front-end}. O objetivo imediato e permitir uma demonstracao apresentavel em curto prazo, com escopo controlado, responsividade real e navegacao geografica progressiva, sem criar uma implementacao descartavel.

A V1 deve entregar uma unica dashboard epidemiologica contendo:

\begin{itemize}[leftmargin=*]
    \item cabecalho institucional do NSS/UENF;
    \item mapa vetorial interativo como elemento central;
    \item navegacao geografica progressiva Brasil $\rightarrow$ Sudeste $\rightarrow$ Rio de Janeiro;
    \item filtros por doenca, ano e mes;
    \item painel contextual com selecao atual, total do recorte quando semanticamente valido e legenda;
    \item ranking dos municipios com dados no nivel do Rio de Janeiro;
    \item estados de carregamento, erro, ausencia de dados e ausencia de cobertura;
    \item suporte a mock e API real por uma mesma abstracao de acesso a dados;
    \item comportamento mobile-first e adaptacao para telas maiores.
\end{itemize}

Nao fazem parte desta primeira entrega: autenticacao visual, administracao de usuarios, dashboards multiplas, filtros demograficos, cartografia de ruas, geocodificacao, GPS, WebSocket, SSR, exportacoes ou um sistema de design complexo.

\section{Arquitetura do NSS e limite do front-end}

A arquitetura de referencia do NSS separa quatro componentes:

\begin{enumerate}[leftmargin=*]
    \item \textbf{Front-end}: React + TypeScript, responsavel pela interface e visualizacoes;
    \item \textbf{Java / Spring Boot}: API principal, autenticacao, autorizacao, regras de negocio e leitura do PostgreSQL;
    \item \textbf{PostgreSQL}: persistencia dos dados consolidados e dados da aplicacao;
    \item \textbf{Python / PySUS}: pipeline de ingestao e transformacao do SINAN, fora do caminho sincrono de consulta do usuario.
\end{enumerate}

\begin{decision}
O front-end nao acessa PostgreSQL e nao chama o pipeline Python/PySUS. Toda consulta epidemiologica da aplicacao deve passar pela API Java. A geometria cartografica local em GeoJSON e um recurso estatico do front-end e nao representa acesso ao pipeline epidemiologico.
\end{decision}

O fluxo desejado e:

\begin{lstlisting}
Usuario -> Front-end React -> API Java -> PostgreSQL
                                      ^
                                      |
                              Pipeline Python/PySUS
\end{lstlisting}

\section{Estado atual do componente Java/Spring Boot}

O projeto Java analisado ja possui elementos importantes para ocupar o papel de backend principal, mas ainda representa apenas o inicio dessa responsabilidade.

\subsection{O que ja existe}

A base atual utiliza Java 21, Spring Boot 4.1.1, Spring Web MVC, Spring Data JPA, Bean Validation, PostgreSQL, Flyway, Spring Security, DTOs, MapStruct, Lombok, Dockerfile e Docker Compose.

A entidade atualmente implementada e \texttt{User}, com CRUD HTTP em \texttt{/NSS/users}. Ainda nao existem os recursos epidemiologicos que o front-end desta V1 devera consumir.

\subsection{O que ainda falta para o papel arquitetural definitivo}

Ainda precisam ser implementados no backend Java:

\begin{itemize}[leftmargin=*]
    \item entidades e repositories para agregados epidemiologicos;
    \item endpoints para doencas e agregados municipais;
    \item migrations Flyway para o schema;
    \item contrato padronizado de erros;
    \item configuracao CORS para o front-end;
    \item especificacao OpenAPI;
    \item testes de integracao do banco e da API epidemiologica;
    \item autenticacao e autorizacao definitivas em fase posterior.
\end{itemize}

\begin{warningbox}
O arquivo \texttt{src/main/resources/application.yml} do projeto Java precisa ser revisado: no estado analisado, ele contem definicoes equivalentes a Docker Compose em vez de propriedades Spring. A definicao dos servicos pertence ao \texttt{compose.yaml}; o \texttt{application.yml} deve conter configuracoes da aplicacao.
\end{warningbox}

\section{Contrato de dados PostgreSQL para a V1}

A tabela epidemiologica proposta para a primeira integracao e deliberadamente pequena e espelha o Gold mensal que o pipeline atual ja produz.

\begin{lstlisting}[language=SQL]
CREATE TABLE sinan_case_aggregate (
    id            BIGSERIAL PRIMARY KEY,
    disease       VARCHAR(20)  NOT NULL,
    source_year   SMALLINT     NOT NULL,
    year          SMALLINT     NOT NULL CHECK (year >= 2000),
    month         SMALLINT     NOT NULL CHECK (month BETWEEN 1 AND 12),
    cd_uf         CHAR(2)      NOT NULL,
    nm_uf         VARCHAR(100) NOT NULL,
    cd_mun        CHAR(7)      NOT NULL,
    nm_mun        VARCHAR(150) NOT NULL,
    cases_total   INTEGER      NOT NULL CHECK (cases_total >= 0),
    batch_id      VARCHAR(32)  NOT NULL,
    ingested_at   TIMESTAMPTZ  NOT NULL,

    CONSTRAINT uq_sinan_case_aggregate_v1 UNIQUE
      (disease, source_year, year, month, cd_uf, cd_mun)
);

CREATE INDEX idx_sinan_case_filter
ON sinan_case_aggregate (disease, year, month, cd_uf);
\end{lstlisting}

\begin{decision}
A V1 utiliza apenas casos mensais agregados por municipio. Sexo, classificacao final, evolucao, semana epidemiologica, idade e outras dimensoes ficam fora desta primeira demonstracao.
\end{decision}

A cobertura de dados de demonstracao fica limitada a quatro municipios do Rio de Janeiro: Campos dos Goytacazes, Sao Joao da Barra, Itaperuna e Macae.

\section{Contrato minimo da API Java para o front-end}

A API V1 deve usar prefixo versionado \texttt{/api/v1}.

\subsection{Health check}

\begin{lstlisting}
GET /api/v1/health

200 OK
{
  "status": "ok"
}
\end{lstlisting}

\subsection{Doencas disponiveis}

\begin{lstlisting}
GET /api/v1/diseases

200 OK
{
  "items": ["DENG"]
}
\end{lstlisting}

\subsection{Agregado municipal}

\begin{lstlisting}
GET /api/v1/epidemiology/municipalities?disease=DENG&year=2026&month=1

200 OK
{
  "disease": "DENG",
  "year": 2026,
  "month": 1,
  "items": [
    {
      "cdUf": "33",
      "nmUf": "Rio de Janeiro",
      "cdMun": "3301009",
      "nmMun": "Campos dos Goytacazes",
      "casesTotal": 123
    }
  ]
}
\end{lstlisting}

Esse endpoint permite ao front-end calcular ranking municipal, total do recorte municipal disponivel e intensidade visual dos municipios cobertos.

\section{Arquitetura do front-end}

\subsection{Stack}

Para a V1, utilizar:

\begin{itemize}[leftmargin=*]
    \item React;
    \item TypeScript com \texttt{strict} habilitado;
    \item Vite;
    \item Tailwind CSS;
    \item TanStack Query para estado de servidor;
    \item Fetch API nativa como cliente HTTP;
    \item React Simple Maps para renderizacao SVG das geometrias GeoJSON.
\end{itemize}

Evitar Redux, Zustand, Next.js, bibliotecas de formulario, bibliotecas de mapas baseadas em tiles e demais dependencias sem necessidade demonstrada.

\subsection{Estrutura sugerida}

\begin{lstlisting}
src/
|-- app/
|   |-- App.tsx
|   `-- providers.tsx
|-- api/
|   |-- http.ts
|   `-- epidemiology.ts
|-- features/
|   `-- dashboard/
|       |-- DashboardPage.tsx
|       |-- components/
|       |   |-- GeographicMap/
|       |   |   |-- GeographicMap.tsx
|       |   |   |-- RegionLayer.tsx
|       |   |   |-- StateLayer.tsx
|       |   |   |-- MunicipalityLayer.tsx
|       |   |   |-- MapTooltip.tsx
|       |   |   `-- MapLegend.tsx
|       |   |-- GeographyBreadcrumb.tsx
|       |   |-- FilterPanel.tsx
|       |   |-- FilterDrawer.tsx
|       |   `-- MunicipalityRanking.tsx
|       `-- hooks/
|           |-- useMapNavigation.ts
|           `-- useEpidemiologyQuery.ts
|-- data/
|   |-- dataSource.ts
|   |-- apiDataSource.ts
|   `-- mockDataSource.ts
|-- mocks/
|   `-- epidemiology.ts
|-- types/
|   `-- epidemiology.ts
|-- index.css
`-- main.tsx
\end{lstlisting}

\subsection{Abstracao de fonte de dados}

A interface visual nao deve saber se os dados vieram do mock ou do Java.

\begin{lstlisting}
interface EpidemiologyDataSource {
  listDiseases(): Promise<string[]>;
  getMunicipalityCases(filters: {
    disease: string;
    year: number;
    month: number;
  }): Promise<MunicipalityCasesResponse>;
}
\end{lstlisting}

O modo e escolhido por variavel de ambiente:

\begin{lstlisting}
VITE_API_BASE_URL=http://localhost:8080
VITE_USE_MOCKS=true
\end{lstlisting}

A troca para a API real nao deve exigir alteracao nos componentes visuais.

\section{Arquitetura cartografica da V1}

\subsection{Decisao: GeoJSON vetorial em vez de PNGs em grid}

A cartografia deve ser gerada previamente por Python usando \texttt{geobr}/GeoPandas, mas persistida no front-end como GeoJSON simplificado. Nao utilizar imagens PNG individuais remontadas em grid como base da interacao.

Motivos:

\begin{itemize}[leftmargin=*]
    \item SVG/GeoJSON preserva as geometrias e permite hitboxes reais;
    \item hover, focus, clique e estilos ficam associados ao poligono correto;
    \item a responsividade nao depende de alinhar imagens rasterizadas;
    \item cores de heatmap podem ser aplicadas dinamicamente;
    \item a mesma geometria funciona em diferentes resolucoes de tela;
    \item acessibilidade por teclado e mais simples de implementar;
    \item nao e necessario gerar milhares de imagens de municipios.
\end{itemize}

\subsection{Arquivos cartograficos necessarios}

A V1 precisa somente destes recursos estaticos:

\begin{lstlisting}
public/maps/
|-- brazil-regions.geojson
|-- southeast-states.geojson
`-- rj-municipalities.geojson
\end{lstlisting}

O primeiro deve conter as cinco regioes brasileiras; o segundo, somente ES, MG, RJ e SP; o terceiro, os municipios do estado do Rio de Janeiro.

\subsection{Script de geracao}

Criar um script de desenvolvimento, por exemplo \texttt{scripts/generate\_maps.py}, responsavel por:

\begin{enumerate}[leftmargin=*]
    \item ler as regioes com \texttt{geobr.read\_region};
    \item ler os estados e filtrar o Sudeste;
    \item ler os municipios do RJ com \texttt{geobr.read\_municipality(code\_muni="RJ")};
    \item manter os codigos oficiais necessários em \texttt{properties};
    \item simplificar a geometria somente se necessario para reduzir tamanho, sem deformacao perceptivel;
    \item exportar GeoJSON valido em \texttt{public/maps};
    \item nao depender de Python em runtime no navegador.
\end{enumerate}

\begin{decision}
O Python desta secao e apenas uma ferramenta de build cartografico. Ele nao substitui o pipeline PySUS e nao entra no caminho de consulta da aplicacao.
\end{decision}

\section{Modelo de navegacao geografica}

O mapa deve funcionar como uma navegacao hierarquica de tres niveis na V1:

\begin{lstlisting}
BRASIL
  |- Norte
  |- Nordeste
  |- Centro-Oeste
  |- Sul
  `- Sudeste ----------> SUDESTE
                           |- ES
                           |- MG
                           |- SP
                           `- RJ -------> RIO DE JANEIRO
                                          |- Campos dos Goytacazes
                                          |- Sao Joao da Barra
                                          |- Itaperuna
                                          `- Macae
\end{lstlisting}

\subsection{Nivel 1: Brasil por regioes}

No carregamento inicial, exibir as cinco regioes brasileiras.

Todas devem possuir hover, focus e selecao visual. Entretanto, somente \textbf{Sudeste} permite drill-down na V1. As demais regioes podem ser selecionadas para mostrar nome e estado de cobertura, mas nao devem navegar para outro nivel.

Nao apresentar \texttt{0 casos} para regioes sem cobertura. Utilizar texto como \texttt{Dados ainda nao disponiveis nesta V1}.

\subsection{Nivel 2: estados do Sudeste}

Ao ativar Sudeste, o mapa deve mudar para ES, MG, RJ e SP. Todos os estados devem ser identificaveis por hover/focus/toque.

Somente o \textbf{Rio de Janeiro} permite drill-down na V1. SP, MG e ES devem informar ausencia de cobertura desta demonstracao, sem inventar total zero.

\subsection{Nivel 3: municipios do Rio de Janeiro}

Ao ativar RJ, renderizar a malha municipal do estado para preservar a forma geografica completa.

Todos os municipios podem aparecer graficamente, mas somente quatro possuem dados interativos epidemiologicos nesta V1:

\begin{itemize}[leftmargin=*]
    \item Campos dos Goytacazes;
    \item Sao Joao da Barra;
    \item Itaperuna;
    \item Macae.
\end{itemize}

Os municipios sem cobertura devem permanecer visualmente neutros e nao podem ser interpretados como municipios com zero casos.

\subsection{Breadcrumb e retorno}

Exibir breadcrumb do nivel atual:

\begin{lstlisting}
Brasil
Brasil > Sudeste
Brasil > Sudeste > Rio de Janeiro
\end{lstlisting}

Os niveis anteriores devem ser clicaveis. Tambem deve existir uma acao clara \texttt{Voltar} em contexto apropriado.

\section{Estado de navegacao e estado epidemiologico}

Separar explicitamente tres conceitos no codigo.

\subsection{Nivel do mapa}

\begin{lstlisting}
type MapLevel =
  | "BRAZIL_REGIONS"
  | "SOUTHEAST_STATES"
  | "RJ_MUNICIPALITIES";
\end{lstlisting}

\subsection{Geografia selecionada}

\begin{lstlisting}
type GeographySelection = {
  level: "region" | "state" | "municipality";
  code: string;
  name: string;
};
\end{lstlisting}

\subsection{Filtros epidemiologicos}

\begin{lstlisting}
type EpidemiologyFilters = {
  disease: string;
  year: number;
  month: number;
};
\end{lstlisting}

\begin{decision}
\texttt{mapLevel}, \texttt{selectedGeography} e \texttt{epidemiologyFilters} nao devem ser fundidos em um unico estado monolitico. Essa separacao reduz acoplamento e simplifica a evolucao da navegacao.
\end{decision}

\section{Join entre geometria e dados epidemiologicos}

GeoJSON responde \textit{onde} a unidade geografica esta e qual seu codigo. A API responde \textit{quantos casos} existem no filtro selecionado.

A ligacao deve usar codigo oficial, nunca nome textual como chave principal.

\begin{lstlisting}
GeoJSON feature
  properties.code_muni = "3301009"
            |
            v
API item
  cdMun = "3301009"
  casesTotal = 120
\end{lstlisting}

A geometria nao deve armazenar valores epidemiologicos fixos.

\section{Cobertura e semantica de dados}

\subsection{Diferenciar zero de ausencia de dados}

A interface deve distinguir visual e semanticamente:

\begin{itemize}[leftmargin=*]
    \item \textbf{sem cobertura}: nao ha dados disponiveis nesta V1;
    \item \textbf{zero casos}: existe cobertura para a unidade e o total calculado e zero;
    \item \textbf{valor positivo}: existe cobertura e ha casos registrados no recorte.
\end{itemize}

Cinza neutro deve representar \textbf{sem cobertura}; a menor intensidade da escala epidemiologica deve representar zero ou baixa incidencia conforme a regra visual escolhida.

\subsection{Agregacoes incompletas}

\begin{warningbox}
O front-end nao deve somar os quatro municipios cobertos e apresentar esse valor como total do estado do Rio de Janeiro, da regiao Sudeste ou do Brasil. Esses seriam totais incompletos e semanticamente incorretos.
\end{warningbox}

Na V1, niveis sem cobertura completa servem primordialmente como navegacao. Totais agregados so podem ser apresentados quando o backend declarar cobertura suficiente para a unidade geografica.

\subsection{Metadados de demonstracao}

O mock pode declarar capacidades de drill-down separadamente da disponibilidade de dados:

\begin{lstlisting}
const demoCoverage = {
  drilldownEnabled: {
    regions: ["SE"],
    states: ["RJ"],
  },
  caseDataAvailableForMunicipalities: [
    "3301009",
    "3305000",
    "3302205",
    "3302403",
  ],
};
\end{lstlisting}

\section{Especificacao visual e responsividade}

\subsection{Principio mobile-first}

A interface deve ser implementada primeiro para viewport estreita e depois expandida progressivamente. Nao criar um desktop fixo para depois apenas empilhar elementos no celular.

\subsection{Mobile}

Em telas pequenas:

\begin{itemize}[leftmargin=*]
    \item o mapa ocupa \textbf{100\% da largura disponivel};
    \item o painel de filtros e informacoes fica oculto por padrao;
    \item uma acao visivel abre o painel como Drawer/Sheet lateral ou sobreposto;
    \item o Drawer deve fechar por botao, toque fora e tecla \texttt{Escape};
    \item ao abrir, o foco deve ir para o painel e permanecer controlado enquanto estiver ativo;
    \item nao depender de hover para nenhuma informacao essencial;
    \item toque em uma geografia deve equivaler a clique/selecionar.
\end{itemize}

Representacao conceitual:

\begin{lstlisting}
+------------------------------+
| NSS / UENF                   |
+------------------------------+
| breadcrumb                   |
+------------------------------+
|                              |
|            MAPA              |
|           100%               |
|                              |
+------------------------------+
| [Filtros e informacoes]      |
+------------------------------+
\end{lstlisting}

\subsection{Telas medias e grandes}

Quando houver largura suficiente, trocar para duas colunas:

\begin{lstlisting}
+--------------------------------------+------------------+
|                                      |                  |
|                MAPA                  | Filtros          |
|              ~ 70%                   | Selecao atual    |
|                                      | Informacoes      |
|                                      | Legenda          |
|                                      |                  |
+--------------------------------------+------------------+
\end{lstlisting}

A proporcao inicial deve ser aproximadamente 70/30, mas o breakpoint e a distribuicao podem ser ajustados apos teste visual. Nao assumir obrigatoriamente \texttt{md}; testar \texttt{md} e \texttt{lg} e usar o primeiro ponto em que a leitura do painel e do mapa permaneça confortavel.

O painel lateral deve preferir \texttt{position: sticky} dentro de CSS Grid/Flex, e nao \texttt{position: fixed}, para permanecer integrado ao fluxo da pagina.

\subsection{Conteudo do painel}

O painel deve conter, no maximo:

\begin{enumerate}[leftmargin=*]
    \item filtro de doenca;
    \item filtro de ano;
    \item filtro de mes;
    \item nome da geografia selecionada;
    \item informacao de cobertura;
    \item total apenas quando semanticamente valido;
    \item legenda do mapa;
    \item indicador discreto \texttt{DEMO} quando estiver usando mock.
\end{enumerate}

\section{Interacao: mouse, toque e teclado}

Toda geografia interativa deve ter comportamentos equivalentes para diferentes dispositivos.

\begin{tabularx}{\textwidth}{@{}L{3.5cm}Y@{}}
\toprule
\textbf{Entrada} & \textbf{Comportamento esperado} \\
\midrule
Mouse hover & realcar poligono e mostrar nome/informacao curta \\
Focus por teclado & mesmo destaque essencial do hover \\
Clique & selecionar geografia e, quando habilitado, fazer drill-down \\
Toque & selecionar; se a unidade permitir drill-down, executar navegacao de forma previsivel \\
Enter/Space & equivalente a ativacao por clique para geografia focada \\
Escape & fechar Drawer/Sheet quando aberto \\
\bottomrule
\end{tabularx}

Nao esconder informacao indispensavel exclusivamente em tooltip de hover.

\section{Visualizacao epidemiologica no mapa}

A intensidade de cor deve ser derivada de \texttt{casesTotal}. A V1 pode utilizar uma escala discreta simples, desde que a legenda explique seu significado.

Regras obrigatorias:

\begin{itemize}[leftmargin=*]
    \item municipio sem cobertura: cinza neutro;
    \item municipio com cobertura e zero casos: estilo valido distinto de sem cobertura;
    \item municipio com valor positivo: intensidade crescente;
    \item estado/regiao sem cobertura completa: nao receber cor que sugira total epidemiologico completo;
    \item contraste suficiente entre estados normal, hover/focus e selecionado.
\end{itemize}

Nao e necessario implementar uma escala epidemiologica sofisticada ou normalizada por populacao nesta V1.

\section{Zoom e transicoes}

A V1 nao precisa calcular zoom GIS sofisticado nem animar bounds dinamicamente. Podem ser utilizados presets de visualizacao para os tres niveis.

\begin{lstlisting}
const MAP_VIEWS = {
  brazil: { center: [-54, -15], zoom: 1 },
  southeast: { center: [-45, -20], zoom: 2.2 },
  rioDeJaneiro: { center: [-42.5, -22.2], zoom: 5 },
};
\end{lstlisting}

Os valores finais devem ser ajustados visualmente. Transicoes curtas podem ser usadas, mas nao sao criterio de aceite se colocarem o prazo em risco.

\section{Tela unica da V1}

A pagina deve seguir esta hierarquia visual:

\begin{enumerate}[leftmargin=*]
    \item Header institucional;
    \item breadcrumb geografico;
    \item area principal contendo mapa + painel responsivo;
    \item ranking municipal abaixo, exibido quando fizer sentido no nivel do RJ;
    \item estados de loading, error, empty e sem cobertura integrados ao fluxo.
\end{enumerate}

Os antigos cards KPI no topo deixam de ser obrigatorios. Informacoes resumidas devem preferencialmente aparecer no painel contextual para preservar o mapa como elemento principal.

\section{Mock da demonstracao}

O mock deve usar o mesmo formato da futura resposta Java.

\begin{lstlisting}
DENG / 2026 / janeiro
- Campos dos Goytacazes: 120
- Macae: 85
- Itaperuna: 42
- Sao Joao da Barra: 18
\end{lstlisting}

Os valores sao sinteticos e devem ser identificados como demonstracao. Nunca apresenta-los como estatistica epidemiologica real.

\section{Estados obrigatorios da interface}

Toda consulta epidemiologica deve tratar:

\begin{itemize}[leftmargin=*]
    \item \textbf{loading}: skeleton/placeholder sem deslocamento importante de layout;
    \item \textbf{error}: mensagem curta e acao de tentar novamente;
    \item \textbf{empty}: existe cobertura, mas nao ha registros para o filtro;
    \item \textbf{no coverage}: a V1 nao possui cobertura para a unidade geografica;
    \item \textbf{success}: renderizacao normal de dados.
\end{itemize}

\section{Escopo explicitamente fora da V1}

Para proteger a entrega, nao implementar agora:

\begin{itemize}[leftmargin=*]
    \item login, cadastro, recuperacao de senha ou RBAC;
    \item tela CRUD de usuarios;
    \item consumo direto do FastAPI antigo;
    \item acesso direto ao PostgreSQL;
    \item filtros por sexo, classificacao, evolucao, idade ou semana epidemiologica;
    \item todos os municipios do Brasil como areas epidemiologicas ativas;
    \item drill-down fora de Sudeste e RJ;
    \item mapas de ruas, satelite ou tiles;
    \item exportacao PDF/CSV;
    \item notificacoes em tempo real;
    \item WebSocket;
    \item SSR/Next.js;
    \item dashboard builder;
    \item cobertura de testes extensa;
    \item animacao cartografica sofisticada.
\end{itemize}

\section{Ordem de implementacao recomendada ao Codex}

\begin{enumerate}[leftmargin=*]
    \item inicializar Vite + React + TypeScript;
    \item configurar Tailwind e TanStack Query;
    \item criar tipos epidemiologicos e estados de navegacao geografica;
    \item criar \texttt{EpidemiologyDataSource};
    \item implementar \texttt{mockDataSource};
    \item criar \texttt{scripts/generate\_maps.py};
    \item gerar e versionar os tres GeoJSONs da V1;
    \item implementar o nivel Brasil/regioes;
    \item implementar drill-down do Sudeste;
    \item implementar drill-down do RJ com malha municipal;
    \item fazer join por codigo IBGE entre GeoJSON e mock;
    \item implementar breadcrumb e navegacao de retorno;
    \item implementar painel lateral desktop e Drawer mobile reutilizando o mesmo conteudo de filtros;
    \item implementar legenda e estados sem cobertura;
    \item implementar ranking dos quatro municipios;
    \item implementar \texttt{apiDataSource};
    \item validar responsividade, teclado, toque e mouse;
    \item executar \texttt{npm run build} e \texttt{npm run lint}.
\end{enumerate}

\section{Criterios de aceite da V1}

\begin{tabularx}{\textwidth}{@{}L{5cm}Y@{}}
\toprule
\textbf{Criterio} & \textbf{Aceite} \\
\midrule
Build & \texttt{npm run build} termina com sucesso \\
Lint & sem erros bloqueadores \\
Mobile-first & mapa ocupa 100\% da largura e painel abre como Drawer/Sheet \\
Desktop & mapa e painel aparecem em duas colunas em breakpoint adequado \\
Mapa Brasil & cinco regioes renderizadas e identificaveis \\
Drill-down regional & somente Sudeste avanca para estados \\
Mapa Sudeste & ES, MG, RJ e SP renderizados \\
Drill-down estadual & somente RJ avanca para municipios \\
Mapa RJ & malha municipal visivel \\
Dados municipais & quatro municipios cobertos recebem dados/heatmap \\
Sem cobertura & nunca e apresentado como zero casos \\
Breadcrumb & permite entender e retornar entre os niveis \\
Interacao & mouse, toque e teclado possuem caminhos equivalentes \\
Filtros & doenca, ano e mes alteram os dados exibidos \\
Join & geometria e dados ligados por codigo oficial, nao por nome \\
Dados & mock e API obedecem aos mesmos tipos \\
Estados & loading, error, empty, no coverage e success existem \\
Acoplamento & nenhum componente visual chama \texttt{fetch} diretamente \\
Arquitetura & navegador conversa apenas com Java no modo real \\
Cartografia & GeoJSON local; nenhum PNG-grid como mecanismo principal \\
Entrega & uma unica dashboard coerente e apresentavel \\
\bottomrule
\end{tabularx}

\section{Evolucao posterior}

Depois da demonstracao, a sequencia recomendada e:

\begin{enumerate}[leftmargin=*]
    \item finalizar persistencia Gold do Python no PostgreSQL;
    \item implementar repository/service/controller epidemiologico no Java;
    \item adicionar OpenAPI e estabilizar DTOs;
    \item integrar o front com a API Java real;
    \item ampliar gradualmente cobertura geografica e declarar explicitamente unidades completas;
    \item implementar autenticacao/autorizacao no Java e somente depois experiencia de login no front;
    \item expandir dimensoes epidemiologicas com contrato SQL e API versionados.
\end{enumerate}

\section{Resumo executivo}

A V1 do \texttt{nss-front-end} e uma dashboard cartografica mobile-first. Seu fluxo principal e Brasil por regioes, drill-down para Sudeste, drill-down para Rio de Janeiro e visualizacao dos quatro municipios atualmente cobertos pelo prototipo. O mapa e requisito central, nao opcional.

A cartografia deve ser vetorial, gerada por Python/\texttt{geobr} em desenvolvimento e versionada como GeoJSON. React Simple Maps renderiza as geometrias como SVG e permite interacao sem depender de mapas de ruas ou servicos externos. O estado geografico permanece separado dos filtros epidemiologicos, e os dados sao associados a geometria por codigos oficiais.

A interface distingue explicitamente \textit{zero casos} de \textit{sem cobertura}, evitando apresentar agregacoes incompletas como totais estaduais, regionais ou nacionais. Em telas pequenas, o mapa ocupa toda a largura e os filtros aparecem em Drawer; em telas maiores, mapa e painel ficam lado a lado com proporcao aproximada 70/30 ajustada por teste visual.

Essa abordagem mantem o escopo enxuto para a demonstracao e, ao mesmo tempo, cria uma base coerente para evoluir a cobertura geografica e conectar a API Java definitiva.

\end{document}
